
\documentclass{article}
\usepackage{colortbl}
\usepackage{makecell}
\usepackage{multirow}
\usepackage{supertabular}

\begin{document}

\newcounter{utterance}

\centering \large Interaction Transcript for game `bbh', experiment `formal\_fallacies', episode 160 with Qwen3.5{-}4B.
\vspace{24pt}

{ \footnotesize  \setcounter{utterance}{1}
\setlength{\tabcolsep}{0pt}
\begin{supertabular}{c@{$\;$}|p{.15\linewidth}@{}p{.15\linewidth}p{.15\linewidth}p{.15\linewidth}p{.15\linewidth}p{.15\linewidth}}
    \# & \multicolumn{2}{c}{Player} && \multicolumn{2}{c}{Game Master} \\
    \hline

    \theutterance \stepcounter{utterance}  
    & & & \multicolumn{4}{p{0.6\linewidth}}{
        \cellcolor[rgb]{0.9,0.9,0.9}{
            \makecell[{{p{\linewidth}}}]{
                \texttt{\tiny{[P1$\langle$GM]}}
                \texttt{Distinguish deductively valid arguments from formal fallacies.} \\
\\ 
\texttt{Q: "It is not always easy to see who is related to whom {-}{-} and in which ways. The following argument pertains to this question: To begin with, Lesley is a close friend of Fernando. Moreover, being a close friend of Fernando or a schoolmate of Lowell is sufficient for being a great{-}grandfather of Leroy. It follows that Lesley is a great{-}grandfather of Leroy."} \\
\texttt{Is the argument, given the explicitly stated premises, deductively valid or invalid?} \\
\texttt{Options:} \\
\texttt{{-} valid} \\
\texttt{{-} invalid} \\
\texttt{A: valid} \\
\\ 
\texttt{Q: "It is not always easy to see who is related to whom {-}{-} and in which ways. The following argument pertains to this question: Whoever is not a great{-}grandfather of Clyde is a stepbrother of Brian. Being an ancestor of Dana is sufficient for not being a great{-}grandfather of Clyde. We may conclude: Everyone who is an ancestor of Dana is a stepbrother of Brian, too."} \\
\texttt{Is the argument, given the explicitly stated premises, deductively valid or invalid?} \\
\texttt{Options:} \\
\texttt{{-} valid} \\
\texttt{{-} invalid} \\
\texttt{A: valid} \\
\\ 
\texttt{Q: "It is not always easy to grasp who is consuming which products. The following argument pertains to this question: Every infrequent user of Paul Mitchell shampoo is either a rare consumer of Nioxin shampoo or a loyal buyer of Caress soap, or both. No regular consumer of Lush soap is a rare consumer of Nioxin shampoo and, in the same time, a loyal buyer of Caress soap. It follows that whoever is an infrequent user of Paul Mitchell shampoo is not a regular consumer of Lush soap."} \\
\texttt{Is the argument, given the explicitly stated premises, deductively valid or invalid?} \\
\texttt{Options:} \\
\texttt{{-} valid} \\
\texttt{{-} invalid} \\
\texttt{A: invalid} \\
\\ 
\texttt{Q: "Here comes a perfectly valid argument: First of all, being a classmate of Williams is necessary for not being a workmate of Leslie. Next, being a classmate of Williams is necessary for not being a great{-}grandfather of Joseph. Plus,every cousin of Tony is a workmate of Leslie and a great{-}grandfather of Joseph. Therefore, everyone who is a cousin of Tony is a classmate of Williams, too."} \\
\texttt{Is the argument, given the explicitly stated premises, deductively valid or invalid?} \\
\texttt{Options:} \\
\texttt{{-} valid } \\
\texttt{{-} invalid} \\
\texttt{A:} \\
            }
        }
    }
     \\ \\

    \theutterance \stepcounter{utterance}  
    & & \multicolumn{4}{p{0.6\linewidth}}{
        \cellcolor[rgb]{0.95,0.95,0.95}{
            \makecell[{{p{\linewidth}}}]{
                \texttt{\tiny{[P1$\leftrightarrow$P]}}
                \texttt{{[}DRAFT{]}} \\
\texttt{To determine if the argument is valid, let's break down the premises and the conclusion using logical notation:} \\
\\ 
\texttt{**Premise 1:** Being a classmate of Williams (\$CW\$) is necessary for not being a workmate of Leslie (\$\textbackslash{}neg WL\$).} \\
\texttt{*   Logical form: \$\textbackslash{}neg WL \textbackslash{}rightarrow CW\$ (If you are not a workmate, you must be a classmate).} \\
\\ 
\texttt{**Premise 2:** Being a classmate of Williams (\$CW\$) is necessary for not being a great{-}grandfather of Joseph (\$\textbackslash{}neg GJ\$).} \\
\texttt{*   Logical form: \$\textbackslash{}neg GJ \textbackslash{}rightarrow CW\$ (If you are not a great{-}grandfather, you must be a classmate).} \\
\\ 
\texttt{**Premise 3:** Every cousin of Tony (\$CT\$) is a workmate of Leslie (\$WL\$) and a great{-}grandfather of Joseph (\$GJ\$).} \\
\texttt{*   Logical form: \$CT \textbackslash{}rightarrow (WL \textbackslash{}land GJ)\$.} \\
\\ 
\texttt{**Conclusion:** Everyone who is a cousin of Tony (\$CT\$) is a classmate of Williams (\$CW\$).} \\
\texttt{*   Logical form: \$CT \textbackslash{}rightarrow CW\$.} \\
\\ 
\texttt{**Analysis:**} \\
\texttt{1.  From Premise 3, we know that anyone who is a cousin of Tony is a **workmate** (\$WL\$) and a **great{-}grandfather** (\$GJ\$).} \\
\texttt{2.  The condition for being a classmate (\$CW\$) is stated in Premises 1 and 2 as being necessary for *not* being a workmate (\$\textbackslash{}neg WL\$) or *not* being a great{-}grandfather (\$\textbackslash{}neg GJ\$).} \\
\texttt{3.  Since the cousin of Tony *is* a workmate and *is* a great{-}grandfather, they do **not** satisfy the conditions (\$\textbackslash{}neg WL\$ or \$\textbackslash{}neg GJ\$) required to trigger the implication that leads to being a classmate.} \\
\texttt{4.  In fact, the premises do not establish any link between being a workmate/great{-}grandfather and being a classmate. One could be a workmate and a great{-}grandfather without being a classmate.} \\
\texttt{5.  Therefore, the conclusion does not logically follow from the premises. The argument commits a formal fallacy (specifically, it confuses necessary conditions with sufficient conditions or fails to link the subject to the predicate).} \\
\\ 
\texttt{The argument is **invalid**.} \\
\\ 
\texttt{A: invalid} \\
            }
        }
    }
    & \\ \\

    \theutterance \stepcounter{utterance}  
    & & \multicolumn{4}{p{0.6\linewidth}}{
        \cellcolor[rgb]{0.95,0.95,0.95}{
            \makecell[{{p{\linewidth}}}]{
                \texttt{\tiny{[P1$\leftrightarrow$P]}}
                \texttt{{[}SELF{-}FEEDBACK{]}} \\
\texttt{FORMAT: ok} \\
\texttt{RULES: ok} \\
\texttt{STRATEGY: ok} \\
\texttt{REASON: The answer correctly identifies the argument as invalid by properly analyzing the necessary vs sufficient condition fallacy.} \\
            }
        }
    }
    & \\ \\

    \theutterance \stepcounter{utterance}  
    & \multicolumn{4}{p{0.6\linewidth}}{
        \cellcolor[rgb]{0.561,0.737,0.561}{
            \makecell[{{p{\linewidth}}}]{
                \texttt{\tiny{[P1$\rangle$GM]}}
                \texttt{A: invalid} \\
            }
        }
    }
    & & \\ \\

    \theutterance \stepcounter{utterance}  
    & & & \multicolumn{2}{p{0.3\linewidth}}{
        \cellcolor[rgb]{0.9,0.9,0.9}{
            \makecell[{{p{\linewidth}}}]{
                \texttt{\tiny{[GM$|$GM]}}
                \texttt{Target: invalid} \\
            }
        }
    }
    & & \\ \\

    \theutterance \stepcounter{utterance}  
    & & & \multicolumn{2}{p{0.3\linewidth}}{
        \cellcolor[rgb]{0.9,0.9,0.9}{
            \makecell[{{p{\linewidth}}}]{
                \texttt{\tiny{[GM$|$GM]}}
                \texttt{game\_result = LOSE} \\
            }
        }
    }
    & & \\ \\

\end{supertabular}
}

\end{document}
